\startchapter[title={Detalles de los modelos lineales}, reference=Anexo_Modelos]

Este anexo reúne dos desarrollos sobre los modelos lineales del
Capítulo~\in[Capitulo_3] que no son necesarios para seguir el texto
principal.

%----------------------------------------------------------------------------------------
\startsection[title={Nivel de servicio con espera voluntaria}, reference=sec:anexo_sla]

Con la notación de la Sección~\in[sec:espera_voluntaria], una consecuencia
de la condición de la Ecuación~\in[eq:espera_condicion] es que, cuando la
espera conviene, la proporción de órdenes que llegan a tiempo queda
determinada por los pesos del objetivo. Con $W^{\text{pen}} = W^{e} = W^{t} = 1$ y sin sucesores que
desplazar, un lote se demora mientras $n^{e}_b - n^{t}_b > W^{c}$, de modo
que se detiene cuando $n^{e}_b - n^{t}_b \approx W^{c}$. Como
$n^{e}_b + n^{t}_b = |\mathcal{N}_b|$, si las fechas de entrega de sus órdenes
están suficientemente repartidas, la proporción de órdenes a tiempo resulta
aproximadamente
\startformula
	\frac{|\mathcal{N}_b| + W^{c}}{2\,|\mathcal{N}_b|},
\stopformula
que depende solo del tamaño del lote y de $W^{c}$ (para $W^{c} < |\mathcal{N}_b|$;
en caso contrario no hay espera).

\stopsection

%----------------------------------------------------------------------------------------
\startsection[title={Validez de $M^{\text{seq}}$ con espera voluntaria}, reference=sec:anexo_mseq]

% TODO: checkear
La cota de
$M^{\text{seq}}$ (Sección~\in[sec:notacion]) se basa en el tiempo total de
trabajo, lo que alcanza cuando cada lote termina en cuanto completa su
ejecución. Con espera voluntaria, en cambio, un lote puede figurar como
terminado mucho después, cerca de la fecha de entrega de sus órdenes, y la
separación entre las finalizaciones de dos lotes de un mismo {\em picker}
puede superar la cota. En ese caso, la restricción de secuencia que debería
estar inactiva pasa a restringir y puede excluir la solución óptima.

Un ejemplo pequeño lo muestra. Sea un único {\em picker} con dos lotes de
duración 10: el primero tiene una orden con fecha de entrega 10, el segundo
tiene cinco órdenes con fecha 2000, y todos los pesos valen 1. El óptimo vale
2010: el primer lote termina en 10 y el segundo se declara terminado en 2000,
sin adelanto ni demora. Con $M^{\text{seq}} = 500$, en cambio, las dos
finalizaciones no pueden separarse más de 490 unidades, y la mejor solución que
admite el modelo vale 5010, con el primer lote demorado hasta 1510.

Para que la cota sea válida también con espera, debe cubrir la fecha de
entrega más lejana. Un lote que termina después de esa fecha tiene todas sus
órdenes atrasadas, y demorarlo solo empeora el objetivo. Por eso siempre
existe una solución óptima en la que, pasada esa fecha, cada {\em picker}
ejecuta sus lotes restantes sin pausas, y ningún lote termina después de
${\max_{n} T^{\text{due}}_n + \sum_{b} \overline{T}_b}$, donde
$\overline{T}_b = F_b + R^{\max} / \min_{p} V_p$ es el tiempo de ejecución de
peor caso del lote $b$, el mismo que usa la cota de la
Sección~\in[sec:notacion]. Sumándole la duración de un lote se obtiene una
cota válida en todos los casos:
\startformula
	M^{\text{seq}} \geq \max_{n \in \mathcal{N}} T^{\text{due}}_n
	+ \sum_{b \in \mathcal{B}} \overline{T}_b + \max_{b \in \mathcal{B}} \overline{T}_b.
\stopformula
Cuando la espera no conviene (Sección~\in[sec:espera_voluntaria]), existe una
solución óptima sin espera y la cota original alcanza.

\stopsection
\stopchapter
